\subsection{Stochastic optimization}
\label{subsec:sgd}

In many statistical learning problems we minimize
\begin{equation}
  f(\theta) = \frac{1}{n}\sum_{i=1}^n f_i(\theta),
  \label{eq:finite-sum}
\end{equation}
e.g.\ each $f_i$ is the loss contributed by data point $i$.
When $n$ is large, evaluating $\nabla f(\theta)$ can be expensive.
The basic idea of SGD is to replace the full gradient $\nabla f(\theta)$ with the gradient of a single randomly sampled term $f_{i_t}(\theta)$ (or a small mini-batch), which makes each iteration cheap.
The running logistic regression example from \cref{ex:logistic-reg} fits this template after normalizing by $n$.
If we define
\[
  f_i(\theta)
  \coloneqq
  \log\bigl(1+e^{x_i^\top\theta}\bigr) - y_i x_i^\top\theta + \frac{m}{2n}\norm{\theta}_2^2,
\]
then $\frac{1}{n}s(\theta)=\frac{1}{n}\sum_{i=1}^n f_i(\theta)$.
So one stochastic gradient step uses the loss contribution from a single observation together with the deterministic ridge gradient $(m/n)\theta$.

\subsubsection{Stochastic gradients and mini-batches}

\begin{algorithm}
  \caption{Stochastic gradient descent (SGD)}
  \label{alg:sgd}
  \begin{algorithmic}[1]
    \Require Functions $f_1,\dots,f_n$, step sizes $(\gamma_t)_{t\ge 0}$, initial point $\theta^{(0)}$, number of steps $T$
    \Ensure Iterate $\theta^{(T)}$
    \For{$t=0,1,\dots,T-1$}
      \State Sample $i_t \sim \mathrm{Unif}\{1,\dots,n\}$
      \State $g_t \gets \nabla f_{i_t}(\theta^{(t)})$
      \State $\theta^{(t+1)} \gets \theta^{(t)} - \gamma_t g_t$
    \EndFor
    \State \Return $\theta^{(T)}$
  \end{algorithmic}
\end{algorithm}

Let $i_t\sim \mathrm{Unif}\{1,\dots,n\}$ and write
\[
g_t = \nabla f_{i_t}(\theta^{(t)}).
\]
Then $g_t$ is an unbiased estimator of the full gradient:
\[
\E[g_t \mid \theta^{(t)}] = \nabla f(\theta^{(t)}).
\]
It is convenient to separate $g_t$ into ``signal'' and ``noise'' by defining the gradient noise
\[
  \xi_t = g_t - \nabla f(\theta^{(t)}),
  \qquad \E[\xi_t \mid \theta^{(t)}]=0.
\]
The conditional covariance of this noise can be written explicitly as
\begin{equation}
  \Sigma(\theta)
  \coloneqq
  \E\bigl[\xi_t\xi_t^\top \mid \theta^{(t)}=\theta\bigr]
  =
  \frac{1}{n}\sum_{i=1}^n \bigl(\nabla f_i(\theta)-\nabla f(\theta)\bigr)\bigl(\nabla f_i(\theta)-\nabla f(\theta)\bigr)^\top.
  \label{eq:sgd-noise-cov}
\end{equation}
This matrix summarizes how heterogeneous the per-sample gradients are at $\theta$.
Even at an optimum $\theta^\star$ (where $\nabla f(\theta^\star)=0$), the per-sample gradients $\nabla f_i(\theta^\star)$ need not vanish, so the stochastic gradient typically has nonzero variance.

In practice one usually averages a batch of $B$ independent indices $i_{t,1},\dots,i_{t,B}$ and uses the mini-batch gradient
\[
g_t^{(B)} = \frac{1}{B}\sum_{j=1}^B \nabla f_{i_{t,j}}(\theta^{(t)}).
\]
This is still unbiased, $\E[g_t^{(B)}\mid \theta^{(t)}]=\nabla f(\theta^{(t)})$.
Moreover, if the indices are sampled independently with replacement, then the noise covariance shrinks by a factor $1/B$:
\begin{equation}
  \E\bigl[(g_t^{(B)}-\nabla f(\theta))(g_t^{(B)}-\nabla f(\theta))^\top \mid \theta^{(t)}=\theta\bigr]
  = \frac{1}{B}\Sigma(\theta).
  \label{eq:minibatch-cov}
\end{equation}
Thus larger batches reduce variance but cost more compute per iteration; mini-batching trades off these two effects.
These considerations already suggest the qualitative behavior of SGD:
with a constant step size the iterates do not settle exactly at $\theta^\star$ (there is persistent noise), while decreasing step sizes and/or averaging can stabilize and improve convergence.

\subsubsection{Step sizes and convergence}
\label{subsec:sgd-stepsizes}

With a fixed step size $\gamma_t\equiv \gamma$, SGD typically does \emph{not} converge exactly to $\theta^\star=\argmin f(\theta)$; it fluctuates around it.
Decreasing step sizes allow convergence to the optimum.
Another standard device is to average the iterates to reduce noise; the bound below is for a weighted average (with weights proportional to the step sizes), which is a convenient form of iterate averaging in theory and practice.

The convergence bound below is standard; see \citet[Thm.~5.3]{garrigos_gower_2024_handbook} for a reference.
We include a proof for completeness.

\begin{theorem}[Averaged SGD bound]
\label{thm:sgd-avg-bound}
Let $\theta^\star = \argmin_{\theta\in\R^d} f(\theta)$ and define
\[
\sigma_\star^2 = \E_i\bigl[\norm{\nabla f_i(\theta^\star) - \nabla f(\theta^\star)}_2^2\bigr]
  = \E_i\bigl[\norm{\nabla f_i(\theta^\star)}_2^2\bigr],
\]
where $i$ is uniform on $\{1,\dots,n\}$ and $\nabla f(\theta^\star)=0$.
Assume:
\begin{itemize}[leftmargin=*]
  \item $f_1,\dots,f_n$ are convex,
  \item $\Hess f_i(\theta) \preceq M_{\max} I$ for all $\theta$ and all $i$.
\end{itemize}
Let $(\theta^{(t)})$ be generated by SGD with step sizes $0<\gamma_t \le 1/(4M_{\max})$.
Define the weighted average iterate
\begin{equation}
  \bar\theta^{\,T} = \frac{\sum_{t=0}^{T-1} \gamma_t \theta^{(t)}}{\sum_{t=0}^{T-1} \gamma_t}.
  \label{eq:theta-bar}
\end{equation}
Then
\begin{equation}
  \E\bigl[f(\bar\theta^{\,T}) - f(\theta^\star)\bigr]
  \le
  \frac{\norm{\theta^{(0)}-\theta^\star}_2^2}{\sum_{t=0}^{T-1}\gamma_t}
  + \frac{2\sigma_\star^2 \sum_{t=0}^{T-1}\gamma_t^2}{\sum_{t=0}^{T-1}\gamma_t}.
  \label{eq:sgd-bound}
\end{equation}
\end{theorem}
\begin{proof}
Fix an iteration $t$ and write $g_t=\nabla f_{i_t}(\theta^{(t)})$.
Expand the squared distance to $\theta^\star$:
\[
\norm{\theta^{(t+1)}-\theta^\star}_2^2
= \norm{\theta^{(t)}-\theta^\star}_2^2 - 2\gamma_t g_t^\top(\theta^{(t)}-\theta^\star) + \gamma_t^2\norm{g_t}_2^2.
\]
Taking conditional expectation given $\theta^{(t)}$ and using unbiasedness $\E[g_t\mid \theta^{(t)}]=\nabla f(\theta^{(t)})$ yields
\begin{equation}
  \E\!\left[\norm{\theta^{(t+1)}-\theta^\star}_2^2 \mid \theta^{(t)}\right]
  = \norm{\theta^{(t)}-\theta^\star}_2^2 - 2\gamma_t \nabla f(\theta^{(t)})^\top(\theta^{(t)}-\theta^\star)
  + \gamma_t^2 \E\!\left[\norm{g_t}_2^2 \mid \theta^{(t)}\right].
  \label{eq:sgd-proof-step1}
\end{equation}
By convexity of $f$, we have $\nabla f(\theta^{(t)})^\top(\theta^{(t)}-\theta^\star) \ge f(\theta^{(t)})-f(\theta^\star)$.

It remains to bound the second moment of $g_t$.
Define
\[
\tilde f_i(\theta) = f_i(\theta) - f_i(\theta^\star) - \nabla f_i(\theta^\star)^\top(\theta-\theta^\star).
\]
Each $\tilde f_i$ is convex and $M_{\max}$-smooth (subtracting an affine function does not change convexity or smoothness), and satisfies $\nabla \tilde f_i(\theta^\star)=0$, so $\theta^\star$ minimizes $\tilde f_i$.
For any convex $M_{\max}$-smooth function $h$ with minimizer $x^\star$ one has
\begin{equation}
  \norm{\nabla h(x)}_2^2 \le 2M_{\max}\bigl(h(x)-h(x^\star)\bigr),
  \label{eq:smooth-grad-bound}
\end{equation}
which follows by applying the smoothness upper bound in \cref{thm:quadratic-bounds} to $h$ at $x$ with the test point $x-\frac{1}{M_{\max}}\nabla h(x)$ and using optimality of $x^\star$.
Applying \eqref{eq:smooth-grad-bound} to $h=\tilde f_i$ gives
\[
\norm{\nabla f_i(\theta)-\nabla f_i(\theta^\star)}_2^2
= \norm{\nabla \tilde f_i(\theta)}_2^2
\le 2M_{\max}\tilde f_i(\theta).
\]
Averaging over $i$ and using $\E_i[\nabla f_i(\theta^\star)]=\nabla f(\theta^\star)=0$ yields
\[
\E_i\norm{\nabla f_i(\theta)-\nabla f_i(\theta^\star)}_2^2
\le 2M_{\max}\bigl(f(\theta)-f(\theta^\star)\bigr).
\]
Finally,
\[
\E_i\norm{\nabla f_i(\theta)}_2^2
\le 2\E_i\norm{\nabla f_i(\theta)-\nabla f_i(\theta^\star)}_2^2 + 2\E_i\norm{\nabla f_i(\theta^\star)}_2^2
\le 4M_{\max}\bigl(f(\theta)-f(\theta^\star)\bigr) + 2\sigma_\star^2.
\]
Plug this into \eqref{eq:sgd-proof-step1} and use $\gamma_t \le 1/(4M_{\max})$ to get
\[
\E\!\left[\norm{\theta^{(t+1)}-\theta^\star}_2^2 \mid \theta^{(t)}\right]
\le \norm{\theta^{(t)}-\theta^\star}_2^2
  - \gamma_t\bigl(f(\theta^{(t)})-f(\theta^\star)\bigr)
  + 2\sigma_\star^2\gamma_t^2.
\]
Rearrange, take full expectation, and sum from $t=0$ to $T-1$ to obtain
\[
\sum_{t=0}^{T-1}\gamma_t\,\E\bigl[f(\theta^{(t)})-f(\theta^\star)\bigr]
\le \norm{\theta^{(0)}-\theta^\star}_2^2 + 2\sigma_\star^2\sum_{t=0}^{T-1}\gamma_t^2.
\]
Since $f$ is convex and $\bar\theta^{\,T}$ is a weighted average, $f(\bar\theta^{\,T}) \le \frac{\sum_t \gamma_t f(\theta^{(t)})}{\sum_t \gamma_t}$.
Dividing by $\sum_{t=0}^{T-1}\gamma_t$ gives \eqref{eq:sgd-bound}.
\end{proof}

The right-hand side of \eqref{eq:sgd-bound} has a simple structure:
the first term is an ``optimization'' term, which decreases when $\sum_t \gamma_t$ is large; the second is a ``noise'' term, which decreases when $\sum_t \gamma_t^2$ is small.
Choosing a step size schedule balances these two effects.

The bound \eqref{eq:sgd-bound} is a clean bias--variance trade-off.
Large step sizes make rapid early progress (large $\sum_t \gamma_t$) but amplify gradient noise (large $\sum_t \gamma_t^2$), leading to a noise floor.
Shrinking step sizes reduces fluctuations and can drive the expected suboptimality to $0$, but typically slows optimization.
At a heuristic level, one can think of SGD as gradient descent perturbed by state-dependent noise whose strength is set by the step size and batch size.
Writing the recursion as
\[
  \theta^{(t+1)} = \theta^{(t)} - \gamma_t \nabla f(\theta^{(t)}) - \gamma_t \xi_t
\]
makes the analogy more concrete.
If $\gamma_t\equiv \gamma$ is small, then on the rescaled time axis $s=t\gamma$ the drift term is precisely the explicit Euler discretization of gradient flow.
Over $O(1/\gamma)$ steps, the accumulated noise is of order $\sqrt{\gamma}$ by a central-limit heuristic, and for a mini-batch of size $B$ this suggests the diffusion approximation
\begin{equation}
  d\Theta_s
  =
  -\nabla f(\Theta_s)\,ds
  + \sqrt{\frac{\gamma}{B}}\,\Sigma(\Theta_s)^{1/2}\,dW_s,
  \label{eq:sgd-sde}
\end{equation}
interpreted informally as a weak-limit statement under additional regularity assumptions; see \citet[\S2]{mandt_hoffman_blei_2015_ct_limit_sgd} for one clean continuous-time derivation in the constant-step regime.
Thus gradient descent is the zero-noise limit, while constant-step SGD is the corresponding Euler--Maruyama discretization with a state-dependent diffusion term.
Near $\theta^\star$, this diffusion does not disappear unless $\gamma$ is reduced or $B$ is increased, which is the dynamical reason constant-step SGD hovers in a neighborhood of the minimizer rather than settling exactly there.
That intuition is consistent with \eqref{eq:sgd-bound}: the same schedule that drives the optimization term down also sets the size of the noise floor.

Some consequences for step size schedules are:
\begin{enumerate}[leftmargin=*,label=(\arabic*)]
  \item \emph{Constant step size.} If $\gamma_t\equiv \gamma$ then
  \[
    \frac{\sum_{t=0}^{T-1}\gamma_t^2}{\sum_{t=0}^{T-1}\gamma_t}=\gamma,
  \]
  so \eqref{eq:sgd-bound} becomes
  \[
    \E[f(\bar\theta^{\,T})-f(\theta^\star)]
    \le \frac{\norm{\theta^{(0)}-\theta^\star}_2^2}{T\gamma} + 2\sigma_\star^2 \gamma.
  \]
  Choosing $\gamma \propto 1/\sqrt{T}$ yields the familiar $O(1/\sqrt{T})$ rate.

  \item \emph{Decaying step sizes.} A common choice is $\gamma_t=\gamma_0/(t+1)^\alpha$ with $1/2<\alpha\le 1$ (often $\alpha=1$), with $\gamma_0$ chosen so that $\gamma_t \le 1/(4M_{\max})$ (or by clipping $\gamma_t$ to satisfy the bound).
  Then $\sum_t \gamma_t = \infty$ but $\sum_t \gamma_t^2 < \infty$, so \eqref{eq:sgd-bound} forces $\E[f(\bar\theta^{\,T})-f(\theta^\star)]\to 0$ as $T\to\infty$.

  \item \emph{Under strong convexity.} If additionally $f$ is strongly convex, then constant-step SGD has a linear transient followed by a noise floor whose size is proportional to $\gamma$; see \citet[Thm.~5.8]{garrigos_gower_2024_handbook}. Decreasing the step size from there, for example with $\gamma_t=\gamma_0/(t+1)$, trades that residual noise for exact convergence, and sharper $O(1/T)$-type bounds are available under stronger assumptions; see \citet[\S5.2]{garrigos_gower_2024_handbook}.
\end{enumerate}
Representative theorems for these schedule-dependent consequences appear in \citet[Thm.~5.3 and \S5.2]{garrigos_gower_2024_handbook}.

\subsubsection{Averaging and practical guidance}
\label{subsubsec:sgd-practical-guidance}

SGD is useful when full gradient computations are bottlenecked by very large data sets or streaming data.
The table below is only a workload heuristic.
It mixes the linear-rate strongly convex regime for GD with the $O(1/\rho)$-type noisy convex scaling suggested by the SGD bounds above, so it should be read as an order-of-magnitude comparison rather than a theorem.
\medskip

A rough trade-off is:
\begin{center}
\begin{tabular}{@{}lccc@{}}
\toprule
& GD & SGD & Mini-batch SGD \\
\midrule
Time / iteration & $n$ & $1$ & $B$ \\
Representative iterations to error scale $\rho$ & $\log(1/\rho)$ & $1/\rho$ & $1/\rho$ \\
Representative time scale & $n\log(1/\rho)$ & $1/\rho$ & $B/\rho$ \\
\bottomrule
\end{tabular}
\end{center}

A few practical implementation notes:
\begin{enumerate}[leftmargin=*,label=(\arabic*)]
  \item Resampling without replacement / shuffling data between epochs is common and can improve performance in practice (also for mini-batches).
  \item Tune learning rates on a small subset first, then scale up.
  A typical decreasing schedule is $\gamma_t = \gamma_0/(t+1)$; another common heuristic is $\gamma_t = \gamma_0/(1+\gamma_0 m t)$ for an estimate of a strong convexity constant $m$.
  \item If possible, monitor both a training-side quantity and a held-out generalization quantity: for example, a full or large-batch estimate of the objective together with validation/test error checked every few epochs.
  Per-mini-batch losses fluctuate too much to support a reliable stopping rule on their own, and watching only training loss can hide the onset of overfitting.
  \item With a constant or piecewise-constant step size, expect a noise-dominated plateau rather than exact convergence.
  Once the monitored metric stops improving for several checks, the right response is usually to reduce $\gamma_t$, increase the batch size, or stop; continuing with the same settings mainly reshuffles the noise floor predicted by \eqref{eq:sgd-bound} and \eqref{eq:sgd-sde}.
  \item If the objective is strongly convex (often true once we add regularization), the last iterate $\theta^{(T)}$ is a reasonable estimator; otherwise use an averaged iterate such as \eqref{eq:theta-bar}, or a \emph{tail average}
  \[
  \bar\theta^{\,T}_{(i)} = \frac{\sum_{t=i}^{T-1} \gamma_t \theta^{(t)}}{\sum_{t=i}^{T-1}\gamma_t}
  \]
  for a sufficiently large burn-in index $i$, chosen after the iterates have entered their final fluctuation regime.
  \item When gradient norms are available only through noise, estimate $\norm{\nabla f(\theta^{(t)})}_2$ on the full data or on a large representative subset before using it as a stopping test.
  This is much more stable than thresholding a single stochastic gradient.
  \item Mini-batching stabilizes noisy gradients; common batch sizes in ML practice are on the order of tens to a few hundreds (e.g.\ $32$--$256$), depending on compute and data.
\end{enumerate}

The stopping logic mirrors the theory.
With a fixed step size $\gamma$ and batch size $B$, the optimization term in \eqref{eq:sgd-bound} decays like $1/(T\gamma)$, while the noise floor stays on the scale of $\gamma/B$.
Once the former is much smaller than the latter, more iterations at the same settings mostly average noise rather than meaningfully decrease the objective.
At that point, either change the schedule or stop and keep an averaged iterate.

Averaging appears constantly in modern ML practice.
One common variant is an exponential moving average of parameters,
\[
\hat\theta^{(t)} = \beta \hat\theta^{(t-1)} + (1-\beta)\theta^{(t)},
\qquad \beta \in (0,1),
\]
which is a geometrically weighted average of past iterates and acts like a low-pass filter that reduces high-frequency noise.
One reason averaging helps is that it reduces the effective gradient noise; \cref{thm:sgd-avg-bound} is one simple guarantee for (weighted) iterate averaging in the convex smooth setting.
EMA is a practically convenient variant that emphasizes recent iterates.
Relatedly, many optimizers used in deep learning maintain exponential moving averages of gradients (and sometimes squared gradients), which can be interpreted as adding smoothing and (diagonal) preconditioning heuristics.

\begin{figure}[t]
  \centering
  \begin{tikzpicture}
  \begin{axis}[
    width=0.9\linewidth,
    height=5.2cm,
    xmin=0, xmax=40,
    ymin=0, ymax=1.2,
    axis lines=left,
    ticks=none,
    xlabel={iteration $t$},
    ylabel={$f(\theta^{(t)})-f(\theta^\star)$},
    clip=false,
  ]
    % A stylized "noisy" objective gap curve + a smoothed (averaged) trend.
    \addplot[stat221Gray!55,thick,domain=0:40,samples=250]
      {0.08 + exp(-0.10*x) + 0.06*sin(18*x) + 0.02*sin(70*x)};
    \addplot[stat221Blue,very thick,domain=0:40,samples=250]
      {0.08 + exp(-0.10*x)};

    \node[stat221Gray!70,anchor=west] at (axis cs:22,0.58) {raw iterates};
    \node[stat221Blue,anchor=west] at (axis cs:22,0.33) {averaged / smoothed};
  \end{axis}
\end{tikzpicture}


  \caption{Typical behavior of SGD on the objective gap: noisy descent with a long tail.
  Averaging (or tail-averaging) stabilizes iterates and makes the underlying decreasing trend easier to see, matching the motivation behind \cref{thm:sgd-avg-bound}.}
  \label{fig:sgd-noisy-gap}
\end{figure}
